\documentclass[a4paper,11pt,fleqn]{scrartcl}%fleqn pour aligner les equations à gauche
\usepackage{../../Styles/Cours}


\newcommand{\chnum}{Chapitre 1}
\newcommand{\chtitle}{Transformation acide-base et pH}
\newcommand{\dtype}{Cours}

\begin{document}
\maketitle

\section{Notions d'acides et de bases}
	\subsection{Définitions}
Selon \textbf{Brönsted} (chimiste danois, 1879 - 1947) :
\begin{itemize*}
	\item un \textbf{acide} est une espèce chimique capable de libérer un ou plusieurs ions hydrogènes \textbf{\ce{H+}}
	\item une \textbf{base}, est une espèce capable de capter un ou plusieurs ions hydrogènes \textbf{\ce{H+}}
\end{itemize*}

\begin{table}[h]
	\begin{tabular}{|p{18cm}}
\textit{Exemple :}\\
- Le sulfure de dihydrogène \ce{H2S} est un acide car il peut libérer un ion hydrogène selon la réaction d'équation :\\
\rule[1cm]{0cm}{0cm}\\
- L'ammoniac \ce{NH3} est une base car elle peut capter un ion hydrogène\\
\rule[1cm]{0cm}{0cm}\\
	\end{tabular}
\end{table}

Certaines espèces possèdent à la fois des propriétés acides et des propriétés basiques : on les appelle des espèces \textbf{amphotères}.
\begin{table}[h]
\begin{tabular}{|m{18cm}}
\textit{Exemple le caractère amphotère de l'eau : }\\
\rule[1.5cm]{0cm}{0cm}\\
\end{tabular}
\end{table}
	\subsection{Couple acide-base}
Un acide et une base forment un \textbf{couple acide base}, noté \textbf{\ce{AH(aq)} / \ce{A-(aq)}}, s'il est possible de passer de l'un à l'autre par un transfert d'un ion hydrogène \ce{H+}. On peut alors écrire une \textbf{demi-équation} de la forme :
$$\ce{AH(aq) <=> A-(aq) + H+}$$
L'acide et la base d'un couple sont dit \textbf{conjugués}.\\
Les \textbf{formules semi-développées} ou les \textbf{schémas de Lewis} des entités acides ou basiques permettent de comprendre leur processus de formation. Une entité acide possède toujours un atome d'hydrogène lié à un atome fortement électronégatif (F, Cl, O, N ...). La liaison entre les deux peut ainsi être facilement rompue.

\begin{table}[h]
\begin{tabular}{|m{18cm}}
\textit{Exemple acide carboxylique : }\\
\rule[1cm]{0cm}{0cm}\\
\end{tabular}
\end{table}

\newpage
Couple acide/base à connaître :
\begin{table}[h]
\centering
\begin{tabular}{|m{6cm}|c|m{7cm}|}
\hline
Eau / ion hydroxyde & \ce{H2O(l)} / \ce{HO-(aq)} &  \rule[1.5cm]{0cm}{0cm} \\
\hline
Ion oxonium / eau & \ce{H3O+(aq)} / \ce{H2O(l)} &  \rule[1.5cm]{0cm}{0cm} \\ 
\hline
Acide carbonique / ion hydrogénocarbonate & \ce{H2CO3 (aq)} / \ce{HCO3- (aq)} & \rule[1.5cm]{0cm}{0cm} \\ 
\hline
Acide carboxylique / ion carboxylate & \ce{RCO2H(aq)} / \ce{RCO2-(aq)} & \rule[1.5cm]{0cm}{0cm}\\
\hline
Ion ammonium / amine & \ce{RNH3+ (aq)} / \ce{RNH2 (aq)} & \rule[1.5cm]{0cm}{0cm}\\
\hline
\end{tabular}
\end{table}
	\subsection{Réaction acide-base}
Une réaction acide-base met en jeux un \textbf{échange d'ions hydrogène $H^{+}$} entre d'acide d'un couple et la base d'un deuxième couple acide/base.\\
\begin{tabular}{|m{\textwidth}}
\textit{Exemple : Réaction de l'acide éthanoïque et de l'eau}\\
1. Écrire la demi-équation associée au couple acide éthanoïque / ion éthanoate. \\
\vspace{2cm}\\
2. Choisir le couple acide/base de l'eau mis en jeu ici, puis la demi-équation associée.\\
\vspace{2cm}\\
3. En déduire l'équation de la réaction \\
\vspace{2cm}\\
\end{tabular}

\section{pH d'une solution}
	\subsection{Définition du pH}
Le \textbf{pH} (potentiel Hydrogène) d'une solution aqueuse est défini par rapport à sa concentration en ions oxonium, notée : [\ce{H3O+}]. Il se détermine selon la relation suivante :
$$pH = -log \frac{[\ce{H3O+}]}{c^0}$$
avec $c^0$ : la concentration standard égale à \SI{1}{\mole\per\liter}.\\
Si la concentration en ion oxonium augmente, alors le pH diminue : la solution devient plus acide.\\

On peut déterminer la concentration en ion oxonium [\ce{H3O+}] à partir du pH d'une solution par la relation :
$$[\ce{H3O+}] = c^0 \cdot 10^{-pH}$$
\end{document}
